\documentclass[11pt]{article}

% ---------- packages ----------
\usepackage[margin=1in]{geometry}
\usepackage{amsmath,amssymb,amsthm}
\usepackage{natbib}
\usepackage{hyperref}
\usepackage{cleveref}
\usepackage{booktabs}
\usepackage{graphicx}
\usepackage{xcolor}

\hypersetup{colorlinks=true,linkcolor=blue!60!black,citecolor=green!50!black,urlcolor=blue!70!black}

\title{Reinforcement Learning, Distillation, and Steering\\of Diffusion Models: A Literature Review}
\author{}
\date{\today}

% ------------------------------------------------------------------
\begin{document}
\maketitle

\begin{abstract}
Diffusion and flow-based generative models have become the dominant paradigm for high-fidelity image, video, and scientific data generation.
A rapidly growing body of work seeks to \emph{control} these models---aligning their outputs with human preferences via reinforcement learning, accelerating sampling through distillation, and steering generation at inference time without retraining.
This review surveys three intertwined threads:
(i) RL-based fine-tuning of diffusion models,
(ii) distillation and few-step generation techniques including flow maps and velocity matching, and
(iii) inference-time scaling and steering via search and reweighting.
We highlight recent advances, identify common theoretical underpinnings, and discuss open problems at the intersection of these areas.
\end{abstract}

\tableofcontents
\newpage

% ==================================================================
\section{Introduction}
\label{sec:intro}

Score-based diffusion models \citep{song2021scorebased,ho2020denoising} and their continuous-time formulation via stochastic differential equations (SDEs) \citep{song2021sde} have achieved remarkable success in image synthesis \citep{dhariwal2021diffusion}, video generation, protein design, and robotic planning.
Despite their quality, two fundamental limitations drive ongoing research:
\begin{enumerate}
    \item \textbf{Alignment.} Pre-trained diffusion models generate samples from the training distribution, which may not match downstream objectives such as aesthetic quality, safety, or task-specific rewards.
    \item \textbf{Efficiency.} Standard samplers require tens to hundreds of neural-function evaluations (NFEs), making real-time deployment costly.
\end{enumerate}

A third, complementary direction---\textbf{inference-time steering}---seeks to guide already-trained models toward high-reward regions without any gradient-based fine-tuning, borrowing ideas from sequential Monte Carlo and search.

This review is organized as follows.
\Cref{sec:evolution} traces the historical evolution of the field.
\Cref{sec:background} reviews the mathematical foundations of diffusion and flow-matching models.
\Cref{sec:rl} surveys reinforcement-learning approaches to diffusion fine-tuning.
\Cref{sec:distillation} covers distillation and few-step generation.
\Cref{sec:steering} discusses inference-time steering and search.
\Cref{sec:applications} highlights applications to planning, manifolds, and scientific domains.
\Cref{sec:lasd} develops a concrete research proposal for linear activation steering of diffusion models.
\Cref{sec:discussion} concludes with open problems.

% ==================================================================
\section{How the Field Evolved: A Ten-Year Arc}
\label{sec:evolution}

The intersection of diffusion models and sequential decision-making has undergone rapid transformation.
We identify five broad phases (\Cref{fig:timeline}).

\begin{figure}[h]
\centering
\small
\setlength{\tabcolsep}{3pt}
\begin{tabular}{@{}r p{0.82\textwidth}@{}}
\toprule
\textbf{Era} & \textbf{Key developments} \\
\midrule
2015--19 & \emph{Theoretical foundations.}
Score matching \citep{vincent2011connection}, deep energy models, and early work on denoising autoencoders lay the mathematical groundwork.
Neural ODEs \citep{chen2018neuralode} introduce continuous-depth networks, planting the seed for flow-based generation. \\[4pt]

2020--21 & \emph{The diffusion revolution.}
DDPM \citep{ho2020denoising} demonstrates that iterative denoising rivals GANs.
Score-based SDEs \citep{song2021scorebased} unify discrete and continuous formulations.
Diffusion models surpass GANs on ImageNet \citep{dhariwal2021diffusion}, triggering an explosion of interest. \\[4pt]

2022 & \emph{Diffusion meets decision-making.}
Janner et al.\ propose \emph{Diffuser} \citep{janner2022planning}---the first work to use a diffusion model as a trajectory planner in offline RL.
Decision Diffuser follows, framing return-conditioned generation as a diffusion process.
These two papers establish the core idea: \emph{a diffusion model can serve as a multimodal policy or planner}, naturally capturing the diversity of optimal behaviors. \\[4pt]

2023 & \emph{Rapid expansion.}
The field broadens along multiple axes simultaneously:
\textbf{(a)}~\emph{Offline RL} becomes the dominant application---diffusion policies replace Gaussian or categorical policies in actor-critic and Q-learning frameworks (IDQL, DiffCPS, MetaDiffuser).
\textbf{(b)}~\emph{Imitation learning} takes off with Diffusion Policy \citep{chi2023diffusion}, which becomes a de facto standard for visuomotor robot learning.
\textbf{(c)}~\emph{Data synthesis}---using diffusion to generate synthetic trajectories for data augmentation (Synthetic Experience Replay, GenAug).
\textbf{(d)}~\emph{Multi-agent and constrained RL} appear (MADiff, AlignDiff).
\textbf{(e)}~The first \emph{survey} paper consolidates the field \citep{zhu2023diffusion}.
Flow matching \citep{lipman2023flow,liu2023flow} emerges as a cleaner training paradigm, and consistency models \citep{song2023consistency} take the first steps toward one-step generation. \\[4pt]

2024 & \emph{Maturation and efficiency.}
Three trends dominate:
\textbf{(a)}~\emph{RL fine-tuning of image diffusion}---DDPO and DPPO \citep{black2024training,fan2024reinforcement} apply PPO to the denoising chain, proving that RLHF works for diffusion models just as it does for LLMs.
\textbf{(b)}~\emph{World models}---diffusion is used to learn environment dynamics (visual details in Atari, autonomous driving).
\textbf{(c)}~\emph{Robustness and safety}---safe offline RL, adversarial perturbation resistance, and preference alignment (PrefPaint, AlignDiff at ICLR).
Efficiency concerns drive work on latent diffusion planning and hierarchical decomposition. \\[4pt]

2025--26 & \emph{Convergence and new frontiers.}
The three threads of this review---RL, distillation, steering---begin to \emph{merge}.
One-step generation matures (MeanFlow, TVM, Drifting), eliminating the sampling bottleneck.
RL moves from the reverse chain to the forward process (DiffusionNFT).
Inference-time steering via SMC (Feynman--Kac) and stochastic flow maps (Meta Flow Maps) offers alignment \emph{without retraining}.
The field expands beyond images: diffusion LLMs with RL (NeurIPS 2025), molecular design, Riemannian manifolds, and long-horizon robotic planning with search (CDGS). \\
\bottomrule
\end{tabular}
\caption{Evolution of diffusion models in RL and controlled generation, 2015--2026.}
\label{fig:timeline}
\end{figure}

\paragraph{The big picture.}
The arc bends from \emph{generation} to \emph{control}.
Early diffusion models were pure samplers---given noise, produce data.
The field then discovered that the same iterative refinement structure makes diffusion models natural \emph{policies} (each denoising step is an action) and \emph{planners} (the trajectory is a sequence of states).
Once this connection was established, the standard RL toolkit---policy gradients, value functions, reward shaping, KL regularization---could be imported wholesale.
In parallel, the efficiency revolution (consistency models $\to$ flow maps $\to$ one-step drifting) removed the computational barrier, while inference-time steering showed that not all control requires gradient-based training.

The current frontier is \emph{unification}: models that are simultaneously fast (few-step), aligned (reward-optimized), and steerable (controllable at test time), with applications spanning robotics, drug design, and language.

% ==================================================================
\section{Background: Diffusion and Flow-Matching Models}
\label{sec:background}

\subsection{Score-Based Diffusion Models}

Diffusion models define a forward noising process that gradually corrupts data $\mathbf{x}_0 \sim p_{\mathrm{data}}$ into Gaussian noise via the SDE
\begin{equation}
    d\mathbf{x}_t = f(\mathbf{x}_t, t)\,dt + g(t)\,d\mathbf{w}_t, \quad t \in [0, T],
\end{equation}
where $f$ is the drift, $g$ the diffusion coefficient, and $\mathbf{w}_t$ a standard Wiener process.
Generation proceeds by simulating the \emph{reverse-time} SDE \citep{anderson1982reverse}:
\begin{equation}
    d\mathbf{x}_t = \bigl[f(\mathbf{x}_t, t) - g(t)^2 \nabla_{\mathbf{x}} \log p_t(\mathbf{x}_t)\bigr]\,dt + g(t)\,d\bar{\mathbf{w}}_t.
\end{equation}
A neural network $s_\theta(\mathbf{x}_t, t) \approx \nabla_{\mathbf{x}} \log p_t(\mathbf{x}_t)$ is trained via denoising score matching \citep{vincent2011connection,song2021scorebased}.
Alternatively, the deterministic probability flow ODE shares the same marginals, enabling exact likelihood computation \citep{song2021sde}.

\subsection{Flow Matching}

Flow matching \citep{lipman2023flow,liu2023flow,albergo2023building} provides a simulation-free training objective for continuous normalizing flows (CNFs).
Given a time-dependent vector field $v_\theta(\mathbf{x}, t)$ that transports a simple prior $p_0$ to the data distribution $p_1$, the flow-matching loss is
\begin{equation}
    \mathcal{L}_{\mathrm{FM}} = \mathbb{E}_{t, \mathbf{x}_t}\bigl[\| v_\theta(\mathbf{x}_t, t) - u_t(\mathbf{x}_t) \|^2\bigr],
\end{equation}
where $u_t$ is a conditional vector field derived from optimal-transport or Gaussian interpolation paths.
Flow matching unifies diffusion and flow-based models and has become the default training paradigm for state-of-the-art systems such as Stable Diffusion 3 and Flux.

\subsection{Connections and Notation}

Throughout this review, we use ``diffusion model'' broadly to include both score-based diffusion and flow-matching models.
The key shared structure is a learned transport map $\mathbf{x}_0 \mapsto \mathbf{x}_T$ (or vice versa) parameterized by a neural velocity or score field, with generation requiring numerical integration of an ODE or SDE.

% ==================================================================
\section{Reinforcement Learning for Diffusion Models}
\label{sec:rl}

\subsection{Motivation}

Pre-trained diffusion models can be viewed as policies in an MDP where each denoising step is an action.
Given a reward function $r(\mathbf{x}_0)$ defined on generated samples, the goal is to fine-tune the model so that it produces high-reward outputs while staying close to the pre-trained distribution (to avoid mode collapse).
This naturally leads to a KL-regularized RL objective:
\begin{equation}
    \max_\theta\; \mathbb{E}_{\mathbf{x}_0 \sim p_\theta}\bigl[r(\mathbf{x}_0)\bigr] - \beta\, \mathrm{KL}\bigl(p_\theta \| p_{\mathrm{pre}}\bigr).
    \label{eq:rl_objective}
\end{equation}

\subsection{Policy-Gradient Methods}

\textbf{DPPO and DDPO.}
\citet{black2024training,fan2024reinforcement} cast the multi-step denoising chain as a finite-horizon MDP and apply proximal policy optimization (PPO).
Each denoising step $\mathbf{x}_{t-1} \sim p_\theta(\cdot | \mathbf{x}_t)$ is treated as an action, with the terminal reward $r(\mathbf{x}_0)$ propagated via standard returns.
These methods demonstrated that RLHF-style training can improve prompt--image alignment, compressibility, and aesthetic scores.

\textbf{FlowGRPO.}
Group Relative Policy Optimization (GRPO) \citep{shao2024deepseekmath} has been adapted for flow-based diffusion models.
It generates a group of samples, computes relative advantages within the group, and updates the policy via a clipped surrogate objective---avoiding the need for a learned value function.

\subsection{DiffusionNFT: RL via the Forward Process}

\citet{zheng2025diffusionnft} propose a fundamentally different approach: rather than discretizing the reverse sampling process, DiffusionNFT optimizes the model \emph{on the forward process via flow matching}.
The method contrasts positive and negative generated samples to establish an improvement direction and incorporates reward signals into a supervised flow-matching objective.
Key advantages include:
\begin{itemize}
    \item Compatibility with arbitrary ODE/SDE solvers (no solver-specific discretization).
    \item No need for likelihood estimation during training.
    \item Up to $25\times$ more efficient than FlowGRPO, achieving a GenEval score of 0.98 vs.\ 0.95 in fewer training steps.
\end{itemize}
This work highlights a growing trend: leveraging the \emph{forward} process structure of flow matching for RL, rather than treating the reverse chain as a black-box MDP.

\subsection{Off-Policy and Value-Based Approaches}

An alternative to on-policy rollouts is to estimate the \emph{value function} at intermediate diffusion states, enabling off-policy updates.
\citet{potaptchik2026meta} (discussed in \Cref{sec:mfm}) show that stochastic one-step posterior sampling can provide efficient value-function estimates, bypassing expensive trajectory simulations.
More broadly, connections to stochastic optimal control \citep{uehara2024finetuning,berner2022optimal} provide a principled framework for diffusion RL, viewing the reverse SDE as a controlled process and the reward as a terminal cost.

% ==================================================================
\section{Distillation and Few-Step Generation}
\label{sec:distillation}

Reducing the number of sampling steps from $\sim\!100$ to $1$--$4$ NFEs is critical for deployment.
We organize this section by the key ideas: consistency/flow maps, velocity matching, and mean-flow approaches.

\subsection{Consistency Models and Flow Maps}

Consistency models \citep{song2023consistency} enforce that all points on the same ODE trajectory map to the same clean output, enabling single-step generation.
\textbf{Flow maps} generalize this idea: rather than mapping to $t=0$, they learn to jump between \emph{arbitrary} pairs of timesteps along the flow.
Self-distillation techniques progressively double the step size, training a student that matches a two-step teacher composed of the same model.

\subsection{Generalised Flow Maps on Riemannian Manifolds}

\citet{davis2025gfm} extend flow maps from Euclidean space to Riemannian manifolds, proposing three self-distillation strategies---Lagrangian, Eulerian, and Progressive.
Their \textbf{Generalised Flow Maps} (GFMs) unify consistency models and shortcut models within a manifold setting, achieving state-of-the-art single- and few-step generation quality on geospatial data, RNA structures, and hyperbolic manifolds.
This work is significant because many scientific data (molecular conformations, climate fields on the sphere) live on non-Euclidean geometries.

\subsection{Terminal Velocity Matching}

\citet{zhou2025tvm} introduce \textbf{Terminal Velocity Matching} (TVM), which models transitions between diffusion timesteps while regularizing at the \emph{endpoint} rather than the starting point.
They prove that TVM provides an upper bound on the 2-Wasserstein distance between data and model distributions when the model is Lipschitz continuous.
Since Diffusion Transformers (DiTs) are not inherently Lipschitz, the authors introduce targeted architectural modifications and a specialized attention kernel for Jacobian--vector product backward passes.
TVM achieves:
\begin{itemize}
    \item \textbf{1-NFE}: 3.29 FID on ImageNet $256 \times 256$.
    \item \textbf{4-NFE}: 1.99 FID, establishing a new state of the art for few-step generation.
\end{itemize}

\subsection{Mean Flows}

\citet{geng2025meanflow} characterize the flow field using the \emph{average velocity}---the integrated velocity from $t$ to $1$---rather than the instantaneous velocity.
They derive an exact relationship between average and instantaneous velocities and train a neural network to predict the average velocity directly.
\textbf{MeanFlow} requires no pre-trained teacher, no distillation, and no curriculum:
\begin{itemize}
    \item Trained from scratch on ImageNet $256 \times 256$, achieving 3.43 FID with a single NFE.
    \item Substantially improves over prior one-step methods and narrows the gap with multi-step models.
\end{itemize}

\subsection{Generative Modeling via Drifting}

\citet{deng2026drifting} propose a novel paradigm where a \emph{drifting field} governs sample movement and reaches equilibrium when the model distribution matches the data distribution.
Unlike diffusion and flow models that require iterative inference, this approach enables one-step generation by evolving the pushforward distribution \emph{during training}.
The method achieves 1.54 FID in latent space on ImageNet $256 \times 256$, demonstrating that the boundary between ``diffusion'' and ``one-step'' models is rapidly dissolving.

\subsection{Summary of Few-Step Methods}

\begin{table}[h]
\centering
\caption{Comparison of few-step generation methods on ImageNet $256 \times 256$.}
\label{tab:fewstep}
\begin{tabular}{lccc}
\toprule
\textbf{Method} & \textbf{NFE} & \textbf{FID $\downarrow$} & \textbf{Key Idea} \\
\midrule
Consistency Models & 1 & 3.55 & ODE trajectory consistency \\
MeanFlow & 1 & 3.43 & Average velocity prediction \\
TVM & 1 & 3.29 & Endpoint regularization \\
Drifting & 1 & 1.54$^*$ & Pushforward evolution \\
TVM & 4 & 1.99 & Endpoint regularization \\
\bottomrule
\multicolumn{4}{l}{\footnotesize $^*$Latent space; not directly comparable to pixel-space FIDs.}
\end{tabular}
\end{table}

% ==================================================================
\section{Inference-Time Steering and Search}
\label{sec:steering}

A complementary approach to RL fine-tuning is to steer a \emph{frozen} model at inference time.
This avoids retraining costs and preserves the base model's diversity, but requires clever use of compute during generation.

\subsection{Feynman--Kac Steering}
\label{sec:fk}

\paragraph{Problem.}
Given a pre-trained diffusion model $p_\theta(\mathbf{x}_0 | c)$ and a reward function $r(\mathbf{x}_0, c)$, we want to sample from the \emph{reward-tilted} distribution
\begin{equation}
    p_{\mathrm{target}}(\mathbf{x}_0 | c) = \frac{1}{Z}\, p_\theta(\mathbf{x}_0 | c)\, \exp\bigl(\lambda\, r(\mathbf{x}_0, c)\bigr),
    \label{eq:tilted}
\end{equation}
which upweights high-reward samples while staying close to the base model.
Fine-tuning methods (\Cref{sec:rl}) modify $\theta$ to approximate this target; steering instead solves the sampling problem \emph{directly at inference time}, leaving $\theta$ frozen.

\paragraph{Approach.}
\citet{singhal2025fksteering} cast this as a \textbf{Feynman--Kac model} over the reverse diffusion trajectory $(\mathbf{x}_T, \mathbf{x}_{T-1}, \ldots, \mathbf{x}_0)$.
They introduce \emph{potential functions} $G_t(\mathbf{x}_t)$ at intermediate timesteps that collectively reconstruct the terminal reward:
\begin{equation}
    \prod_{t=T}^{0} G_t(\mathbf{x}_T, \ldots, \mathbf{x}_t) = \exp\bigl(\lambda\, r(\mathbf{x}_0, c)\bigr).
\end{equation}
The algorithm runs $k$ \emph{particles} (parallel diffusion trajectories) and applies sequential Monte Carlo (SMC):
\begin{enumerate}
    \item \textbf{Propose}: advance each particle one denoising step using the base model.
    \item \textbf{Score}: evaluate the potential $G_t$ for each particle.
    \item \textbf{Resample}: multinomial resampling---particles with high potential scores are duplicated, low-scoring ones are eliminated.
\end{enumerate}
This is strictly more powerful than Best-of-$N$, which only selects at the \emph{end}: FK steering prunes bad trajectories \emph{early}, focusing compute on promising regions of sample space.

\paragraph{Potential design.}
The choice of $G_t$ determines the quality--variance trade-off. Three designs are explored:
\begin{itemize}
    \item \textbf{Difference}: $G_t = \exp\bigl(\lambda[r_\phi(\mathbf{x}_t) - r_\phi(\mathbf{x}_{t+1})]\bigr)$---rewards particles whose intermediate scores are \emph{increasing}.
    \item \textbf{Max}: favors trajectories that reach the highest intermediate reward at any point.
    \item \textbf{Sum}: accumulates rewards across the trajectory.
\end{itemize}
Intermediate rewards $r_\phi(\mathbf{x}_t, c)$ can be computed by evaluating an off-the-shelf reward on the predicted clean sample $\hat{\mathbf{x}}_0(\mathbf{x}_t)$, or by training a dedicated model.

\paragraph{Theoretical guarantee.}
As $k \to \infty$, the empirical particle distribution converges to $p_{\mathrm{target}}$ for \emph{any} valid potential decomposition---the estimator is \textbf{consistent}.

\paragraph{Key results.}
\begin{itemize}
    \item \emph{Text-to-image} (GenEval benchmark): Stable Diffusion v2.1 with $k\!=\!4$ particles outperforms the $3\times$ larger SDXL-DPO fine-tuned model on prompt fidelity. A 0.8B steered model beats a 2.6B fine-tuned model.
    \item \emph{Text diffusion}: gradient-free control of toxicity (0.3\% baseline $\to$ 64.7\% with $k\!=\!8$) and improved linguistic acceptability.
    \item Steering \emph{stacks} with fine-tuning: applying FK steering to an already fine-tuned SDXL-DPO model further improves GenEval from 0.58 to 0.65.
    \item Computational cost is lower than fine-tuning despite multi-particle sampling, since no backward passes or gradient updates are needed.
\end{itemize}
This connects diffusion steering to the broader theme of \emph{inference-time scaling} in language models \citep{snell2024scaling}, suggesting a unified compute-optimal perspective: at some point, spending more compute at inference time yields better returns than spending it on training.

\subsection{Meta Flow Maps for Steering and Alignment}
\label{sec:mfm}

\citet{potaptchik2026meta} introduce \textbf{Meta Flow Maps} (MFMs), extending consistency models into the stochastic domain.
MFMs perform stochastic one-step posterior sampling: from any intermediate noisy state $\mathbf{x}_t$, they generate multiple independent clean-data estimates via a single forward pass.
These estimates enable:
\begin{enumerate}
    \item \textbf{Inference-time steering} without inner rollouts---value functions are estimated by averaging rewards over the one-step samples.
    \item \textbf{Unbiased off-policy fine-tuning}---the differentiable reparameterization trick allows gradient-based reward optimization without on-policy trajectory simulation.
\end{enumerate}
Empirically, the steered MFM sampler outperforms a Best-of-1000 baseline on ImageNet across multiple reward metrics at a fraction of the compute.
MFMs thus bridge the gap between distillation (\Cref{sec:distillation}) and steering, showing that few-step models are not just faster---they also enable qualitatively new control mechanisms.

\subsection{Compositional Diffusion with Guided Search}

\citet{mishra2026cdgs} address long-horizon planning, where compositional diffusion models must combine local sub-task distributions.
When these local distributions are multimodal, naive composition averages incompatible modes, producing infeasible plans.
\textbf{CDGS} integrates \emph{search within the denoising process}:
\begin{itemize}
    \item Population-based sampling explores diverse local mode combinations.
    \item Likelihood-based filtering eliminates infeasible candidates.
    \item Iterative resampling across overlapping segments maintains global consistency.
\end{itemize}
Demonstrated on seven robot manipulation tasks, CDGS also extends to text-guided panoramic image generation and long video synthesis.
This highlights a broader principle: \emph{search and diffusion are complementary}, with search resolving the combinatorial structure that pure denoising cannot.

% ==================================================================
\section{Applications and Extensions}
\label{sec:applications}

\subsection{Diffusion Models for Robotic Planning}

Diffusion models have emerged as powerful trajectory planners in robotics \citep{janner2022planning,chi2023diffusion}.
The key insight is that a diffusion model can represent multi-modal trajectory distributions, capturing diverse strategies for reaching a goal.
RL fine-tuning (\Cref{sec:rl}) and inference-time steering (\Cref{sec:steering}) both apply naturally: the ``reward'' becomes task success or constraint satisfaction, and the compositional structure of CDGS \citep{mishra2026cdgs} addresses the horizon-scaling challenge.

\subsection{Non-Euclidean Domains}

Many scientific applications require generation on non-Euclidean spaces: molecular conformations on $\mathrm{SO}(3)$, climate data on the sphere $S^2$, phylogenetic trees in hyperbolic space.
Generalised Flow Maps \citep{davis2025gfm} extend the distillation toolkit to these settings, while Riemannian score-based models \citep{debortoli2022riemannian} provide the diffusion backbone.
The combination of manifold diffusion with RL and steering remains largely unexplored.

\subsection{Connections to Stochastic Optimal Control}

The RL objective in \Cref{eq:rl_objective} can be reformulated as a stochastic optimal control problem over the reverse SDE \citep{uehara2024finetuning,berner2022optimal}.
The value function satisfies a Hamilton--Jacobi--Bellman (HJB) equation, and the optimal policy tilts the score function by the log-value gradient.
This perspective unifies RL fine-tuning, classifier guidance, and SMC steering as different approximations to the same controlled diffusion.

% ==================================================================
\section{Research Directions: Linear Activation Steering for Diffusion}
\label{sec:lasd}

The preceding sections reveal a gap in the cost--capability landscape of diffusion steering.
RL fine-tuning modifies weights (expensive, permanent).
Feynman--Kac steering requires $k\times$ inference compute (training-free, but costly per generation).
Classifier guidance needs a backward pass per step; CFG doubles the forward passes.
No existing method offers \emph{zero-overhead, training-free, composable} steering of diffusion models.

We propose \textbf{Linear Activation Steering for Diffusion (LASD)}, inspired by the recent discovery that semantic concepts are encoded as linear directions in the residual stream of large transformers \citep{beaglehole2026universal}.
The core idea: extract concept vectors from Diffusion Transformer (DiT) activations using Recursive Feature Machines and steer by vector addition during denoising.

\subsection{Background: Activation Steering in LLMs}

\citet{beaglehole2026universal} demonstrated that for LLMs and VLMs, a concept $\mathcal{C}$ (e.g., ``truthfulness,'' ``formality'') corresponds to a linear direction $\mathbf{v}_\ell \in \mathbb{R}^D$ at each layer $\ell$.
Steering is achieved by modifying activations:
\begin{equation}
    \tilde{\mathbf{A}}_{\ell,i} = \mathbf{A}_{\ell,i} + \varepsilon\, \mathbf{v}_\ell, \quad \forall \text{ token positions } i,
    \label{eq:llm_steering}
\end{equation}
at \emph{zero} additional inference cost (just vector addition).
The concept vectors are extracted via the \textbf{Average Gradient Outer Product (AGOP)} of a kernel ridge regression model trained on (activation, concept-label) pairs: the leading eigenvector of $\mathbf{M}_{\mathrm{AGOP}} = \frac{1}{M}\sum_i \nabla \hat{f}(\mathbf{a}^{(i)})\nabla \hat{f}(\mathbf{a}^{(i)})^\top$ captures the direction of maximal concept variation.
Across 512 concepts on models from Llama-8B to DeepSeek-R1, this achieves 69.1\% steerability, with linear composability ($\mathbf{v} = \sum_k \varepsilon_k \mathbf{v}^{(k)}$) and cross-language transfer.

\paragraph{Prior art in diffusion activation editing.}
The idea of manipulating diffusion model activations is not new.
\citet{kwon2023diffusion} showed that the bottleneck activations of a U-Net diffusion model (the ``h-space'') encode semantic concepts, and that linear arithmetic $\mathbf{h} + \alpha\,\Delta\mathbf{h}$ steers generation.
\citet{gandikota2023concept} found low-rank concept directions in \emph{weight space} via LoRA (Concept Sliders), and \citet{epstein2023selfguidance} used internal attention maps for spatial control (Self-Guidance).
LASD differs from these works in three respects:
(i) it targets the \textbf{DiT architecture} (not U-Nets), which is now the dominant backbone;
(ii) it uses a \textbf{systematic, scalable extraction pipeline} (RFM/AGOP) rather than manual direction finding;
(iii) it explicitly addresses the \textbf{timestep-dependence problem} unique to iterative denoising.
Nevertheless, the contribution is primarily \emph{empirical}---the central bet is whether DiT activations exhibit the same linear concept geometry found in LLMs.

\subsection{Proposed Method}

Let $f_\theta$ be a pre-trained DiT with $L$ layers.
At denoising step $t$ and layer $\ell$, the model computes activations $\mathbf{A}_\ell^{(t)} \in \mathbb{R}^{N \times D}$, where $N$ is the number of spatial tokens (patches) and $D$ the hidden dimension.
LASD steers by:
\begin{equation}
    \tilde{\mathbf{A}}_\ell^{(t)} = \mathbf{A}_\ell^{(t)} + \varepsilon(t)\, \mathbf{v}_\ell^{(t)}\, \mathbf{1}_N^\top,
    \label{eq:lasd}
\end{equation}
where $\mathbf{v}_\ell^{(t)} \in \mathbb{R}^D$ is the concept vector at layer $\ell$ and timestep $t$, $\varepsilon(t)$ is a timestep-dependent steering strength, and $\mathbf{1}_N^\top$ broadcasts across all spatial positions (uniform steering for global concepts).

\paragraph{Concept vector extraction.}
Generate $M$ images from diverse prompts, label each with a concept score $y^{(i)}_\mathcal{C}$ (from a reward model or classifier applied to the \emph{final} image $\mathbf{x}_0$), and record spatially-pooled activations $\mathbf{a}_{\ell,t}^{(i)} = \mathrm{MeanPool}(\mathbf{A}_\ell^{(t)}(\mathbf{x}_t^{(i)}, \mathbf{c}^{(i)})) \in \mathbb{R}^D$.
Apply the RFM pipeline (kernel ridge regression $\to$ AGOP $\to$ leading eigenvector) to the dataset $\{(\mathbf{a}_{\ell,t}^{(i)}, y^{(i)}_\mathcal{C})\}_{i=1}^M$ to obtain $\mathbf{v}_{\ell,t}$.

\paragraph{A note on credit assignment.}
A subtlety: concept labels are defined on $\mathbf{x}_0$ (the final image), but activations are recorded at intermediate steps $t > 0$.
This is analogous to terminal-reward RL---the RFM extracts the direction in activation space that is most predictive of the eventual outcome, effectively learning which intermediate representations ``commit'' to the concept.

\subsection{The Timestep Challenge}
\label{sec:timestep}

This is the central technical question and potential failure point.
In an LLM, the same network processes all tokens in a single forward pass; the residual stream at layer $\ell$ lives in a consistent representational space.
In a diffusion model, the same network is called $T$ times with dramatically different inputs: pure noise at $t = T$, a near-clean image at $t \approx 0$.
The activation manifold shifts continuously with $t$, and a concept direction extracted at one timestep may be meaningless at another.

\paragraph{Diagnostic.}
Before any steering experiments, we propose measuring the \textbf{temporal stability} of concept directions:
\begin{equation}
    S(\mathcal{C}, \ell) = \frac{1}{|\mathcal{T}|^2} \sum_{t_1, t_2 \in \mathcal{T}} \cos\bigl(\mathbf{v}_\ell^{(t_1)},\; \mathbf{v}_\ell^{(t_2)}\bigr),
    \label{eq:stability}
\end{equation}
where $\mathcal{T}$ is a set of sampled timesteps.
If $S \approx 1$, the concept direction is stable and a single time-independent vector suffices.
If $S \approx 0$, the direction rotates across timesteps and per-timestep vectors are necessary.

\paragraph{Three strategies.}
\begin{enumerate}
    \item[\textbf{A.}] \textbf{Shared vector, adaptive scaling.} A single $\mathbf{v}_\ell$ extracted from activations pooled across all $t$, with $\varepsilon(t) = \varepsilon_0 \cdot \sigma(t)^{-\alpha}$ (steer more at low noise, less at high noise). Simplest; serves as baseline.
    \item[\textbf{B.}] \textbf{Time-binned vectors.} Partition the denoising trajectory into $B$ bins (e.g., $B = 4$). Extract a separate $\mathbf{v}_\ell^{(b)}$ per bin. Captures the intuition that coarse layout planning ($t$ large) and fine detail refinement ($t$ small) may encode concepts differently.
    \item[\textbf{C.}] \textbf{Continuous interpolation.} Smooth interpolation between bin vectors via Gaussian-weighted averaging: $\mathbf{v}_\ell(t) = \sum_b w_b(t)\, \mathbf{v}_\ell^{(b)}$, avoiding discontinuities at bin boundaries.
\end{enumerate}
A possible saving grace: DiTs receive the timestep via adaptive LayerNorm (adaLN), which modulates activations based on $t$.
It is \emph{possible} that the concept directions are approximately stable in the adaLN-modulated space, making Strategy A viable.
This is an empirical question.

\subsection{Theoretical Limitations}

Unlike FK steering, which is asymptotically consistent (the particle distribution converges to $p_{\mathrm{target}}$ as $k \to \infty$), LASD has \textbf{no guarantee} that the perturbed denoising trajectory samples from a well-defined target distribution.
Adding $\varepsilon\,\mathbf{v}$ to activations modifies the effective velocity field:
\begin{equation}
    v_{\mathrm{steered}}(\mathbf{x}_t, t) \neq v_\theta(\mathbf{x}_t, t) + \varepsilon \cdot (\text{something simple}),
\end{equation}
because the perturbation propagates nonlinearly through subsequent layers and denoising steps.
The resulting distribution may not correspond to any reward-tilted version of the base model.
At large $\varepsilon$, the trajectory may leave the data manifold entirely, producing artifacts.

This is a fundamental limitation relative to FK steering.
However, the same limitation applies to classifier-free guidance (CFG), which also lacks a rigorous distributional guarantee yet is universally deployed.
The practical question is whether LASD's reward--quality trade-off is favorable, not whether it has an asymptotic theorem.

\subsection{Proposed Experiments}

We outline a staged experimental plan, with clear go/no-go criteria at each stage.

\paragraph{Stage 0: Linearity probing (prerequisite).}
Generate 10,000 images with concept labels.
Record activations at all (layer, timestep) pairs.
Train RFM probes and report classification accuracy as a function of $\ell$ and $t$.
Compare to nonlinear probes (MLP) to quantify the linearity gap.
\textbf{Go/no-go}: if RFM probe accuracy is near chance, the project pivots to a DiT representation analysis paper.

\paragraph{Stage 1: Temporal stability.}
Compute the stability metric $S(\mathcal{C}, \ell)$ (\Cref{eq:stability}) for each concept and layer.
Determine whether Strategy A, B, or C is needed.

\paragraph{Stage 2: Steering evaluation.}
Apply LASD to DiT-XL/2 (class-conditional, ImageNet) and Stable Diffusion 3 (text-to-image).
Evaluate against baselines: CFG, classifier guidance, FK steering ($k = 4, 8$), Concept Sliders, h-space editing adapted to DiT, and a simple mean-difference vector (ablation of the RFM extraction).
Metrics: concept accuracy (reward model), FID (quality preservation), LPIPS diversity, and wall-clock cost ratio.

\paragraph{Stage 3: Composability and scaling.}
Test multi-concept steering ($\mathbf{v} = \sum_k \varepsilon_k \mathbf{v}^{(k)}$) for 2--3 simultaneous concepts.
Measure orthogonality of concept vectors and interference effects.
If resources permit, test on Flux-12B where FK steering ($k = 8$) would be prohibitively expensive.

\paragraph{Stage 4: Ablations.}
Steering strength sweep ($\varepsilon$), layer selection (single layer vs.\ all vs.\ middle-only), spatial mode (uniform vs.\ attention-weighted), number of RFM training samples, and RFM vs.\ simpler extraction methods (PCA, mean-difference, logistic regression).

\subsection{Expected Outcomes and Fallbacks}

\paragraph{Strong success.}
LASD achieves concept accuracy within 10\% of FK steering ($k = 4$) at $1.0\times$ inference cost, with FID degradation under 5\%, across $\geq$50\% of tested concepts.
Multi-concept composition works cleanly.

\paragraph{Moderate success.}
LASD works for global concepts (style, aesthetics) but fails for fine-grained or spatially localized concepts.
Still valuable as a cheap first-pass method---potentially combinable with FK steering (use LASD to bias all $k$ particles, then resample, achieving FK-quality results with $k = 2$ instead of $k = 8$).

\paragraph{Minimum publishable result.}
Clear evidence that DiT activations encode concepts linearly, with a spatiotemporal atlas of concept geometry across layers and timesteps.
This representation analysis would be independently valuable for diffusion model interpretability, even if steering effectiveness is limited.

\paragraph{Failure mode.}
If DiT activations lack linear concept structure (low RFM probe accuracy), the entire approach fails at a foundational level.
The fallback is a mechanistic interpretability study: where and when do DiTs encode semantic information, and why is the geometry different from LLMs?

% ==================================================================
\section{Discussion and Open Problems}
\label{sec:discussion}

\subsection{Theoretical Unification of RL, Distillation, and Steering}

A recurring theme is the convergence of these three threads.
DiffusionNFT \citep{zheng2025diffusionnft} performs RL through the forward (flow-matching) process rather than the reverse chain.
Meta Flow Maps \citep{potaptchik2026meta} use distilled one-step models to enable both steering and off-policy RL.
Feynman--Kac steering \citep{singhal2025fksteering} replaces gradient-based RL with SMC reweighting.
All three optimize variants of the same KL-regularized reward objective (\Cref{eq:tilted}), but through fundamentally different mechanisms---policy gradients on the reverse chain, supervised learning on the forward process, and importance-weighted sampling at test time.
A unified theoretical framework---likely grounded in stochastic optimal control and the HJB equation---that characterizes \emph{when to use which} (as a function of reward structure, compute budget, and alignment precision) remains absent.

\subsection{Inference-Time Compute Scaling Laws}

Language models exhibit clear scaling laws: more chain-of-thought tokens or search budget yields predictable reasoning improvements \citep{snell2024scaling}.
For diffusion models, Feynman--Kac steering and CDGS demonstrate that more particles and search yield better alignment, but there is no analogue of ``scaling laws.''
Key unanswered questions include:
\begin{itemize}
    \item How does reward improvement scale with the number of particles $k$? The SMC consistency guarantee is asymptotic ($k \to \infty$); finite-sample behavior is poorly characterized.
    \item What is the optimal allocation between particles, resampling frequency, and denoising steps for a fixed compute budget?
    \item When is it more efficient to steer at inference time vs.\ fine-tune with RL? The crossover point likely depends on how many generations are needed and how far the target is from the base distribution.
\end{itemize}

\subsection{One-Step Generation: What Do You Lose?}

MeanFlow, TVM, and Drifting have brought one-step FIDs close to multi-step baselines (\Cref{tab:fewstep}).
However, FID measures aggregate distributional similarity and is insensitive to mode dropping, tail coverage, and out-of-distribution robustness.
Compressing a 100-step iterative refinement into a single function evaluation necessarily discards the ``self-correcting'' nature of the denoising process.
Open questions:
\begin{itemize}
    \item Do one-step models lose rare modes? Systematic studies of mode coverage (e.g., via recall metrics or coverage on long-tail categories) are lacking.
    \item How do one-step models interact with steering? FK steering relies on intermediate states for resampling; a one-step model has no intermediates. MFMs partially address this by providing stochastic one-step posterior samples, but the trade-off between speed and steerability is unclear.
\end{itemize}

\subsection{Online RL with Diffusion Policies}

The vast majority of diffusion+RL work is \emph{offline}: the diffusion model is trained on a fixed dataset and fine-tuned without further environment interaction.
Online RL with diffusion policies faces two barriers:
\begin{itemize}
    \item \textbf{Latency.} Multi-step sampling is too slow for real-time control loops. One-step distillation could unlock online deployment, but this has not been demonstrated at scale in closed-loop settings.
    \item \textbf{Exploration.} Diffusion policies naturally represent multi-modal action distributions, which should help exploration, but how to leverage this structure (e.g., Thompson sampling via the stochastic denoising process) is poorly understood.
\end{itemize}

\subsection{Multi-Objective and Constrained Alignment}

Current RL fine-tuning optimizes a single scalar reward.
Real applications require Pareto-optimal trade-offs: image quality vs.\ safety vs.\ diversity, or drug efficacy vs.\ synthesizability vs.\ toxicity.
Constrained diffusion RL (DiffCPS, safe offline RL) and multi-objective formulations are nascent.
Inference-time steering could be particularly natural for multi-objective control---different potential functions for different objectives---but this has not been explored.

\subsection{Reward Hacking and the Alignment Tax}

RL fine-tuning risks \emph{reward hacking}---optimizing a proxy reward at the expense of true quality.
The KL penalty in \Cref{eq:rl_objective} mitigates this, but calibrating $\beta$ is non-trivial: too small and the model over-optimizes the proxy; too large and it barely moves from the base distribution.
Inference-time steering may be more robust since it does not modify model weights and can be ``turned off,'' but it incurs an ``alignment tax'' in the form of additional sampling compute at every generation.
Understanding the robustness--efficiency frontier across methods is an open problem.

\subsection{Beyond Images: Video, 3D, and Science}

Most current work focuses on image generation.
Extending RL fine-tuning and steering to video diffusion models (orders of magnitude more dimensions), 3D generation (geometric constraints), and scientific domains (molecules on $\mathrm{SO}(3)$, climate on $S^2$) presents significant challenges:
\begin{itemize}
    \item \textbf{Reward design.} Image rewards (CLIP, aesthetic scores) are mature; video and 3D rewards are noisy and expensive.
    \item \textbf{Particle methods at scale.} FK steering with $k$ particles on a video model means $k\times$ the memory and compute of an already expensive generation.
    \item \textbf{Manifold RL.} GFMs \citep{davis2025gfm} provide distillation on manifolds, but combining manifold diffusion with RL fine-tuning or steering is essentially untouched.
\end{itemize}

\subsection{Diffusion as Policy vs.\ Diffusion as Planner}

Using a diffusion model as a \emph{policy} (output a single action $\mathbf{a}_t$ given state $\mathbf{s}_t$) vs.\ as a \emph{planner} (output a full trajectory $(\mathbf{s}_0, \mathbf{a}_0, \ldots, \mathbf{s}_H, \mathbf{a}_H)$) yields very different computational and representational trade-offs.
Policies are fast at test time but myopic; planners capture long-horizon structure but scale poorly with horizon length.
Hierarchical approaches (plan at a coarse level, execute with a policy) are a natural compromise, but principled ways to decompose the problem---and how RL fine-tuning should differ across levels---remain open.

% ==================================================================
\section*{Acknowledgments}

We thank Ricky Chen, Stefano Ermon, Qiang Liu, Masatoshi Uehara, Nicholas Boffi, and Michael Albergo for valuable discussions and for making their work publicly available.

% ==================================================================
\bibliographystyle{plainnat}
\begin{thebibliography}{99}

\bibitem[Beaglehole et~al.(2026)]{beaglehole2026universal}
Beaglehole, D., Radhakrishnan, A., Boix-Adser\`{a}, E., and Belkin, M.
\newblock Toward universal steering and monitoring of {AI} models.
\newblock \emph{Science}, 391(6787):787--792, 2026.

\bibitem[Epstein et~al.(2023)]{epstein2023selfguidance}
Epstein, D., Jabri, A., Poole, B., Efros, A.~A., and Holynski, A.
\newblock Diffusion self-guidance for controllable image generation.
\newblock \emph{NeurIPS}, 2023.

\bibitem[Gandikota et~al.(2023)]{gandikota2023concept}
Gandikota, R., Orgad, H., Belinkov, Y., Oh, J., and Materzy\'{n}ska, J.
\newblock Concept sliders: {LoRA} adaptors for precise control in diffusion models.
\newblock \emph{arXiv:2311.12092}, 2023.

\bibitem[Kwon et~al.(2023)]{kwon2023diffusion}
Kwon, M., Jeong, J., and Uh, Y.
\newblock Diffusion models already have a semantic latent space.
\newblock \emph{ICLR}, 2023.

\bibitem[Chen et~al.(2018)]{chen2018neuralode}
Chen, R.~T.~Q., Rubanova, Y., Bettencourt, J., and Duvenaud, D.
\newblock Neural ordinary differential equations.
\newblock \emph{NeurIPS}, 2018.

\bibitem[Zhu et~al.(2023)]{zhu2023diffusion}
Zhu, Z., Zhao, H., He, H., Zhong, Y., Zhang, S., Yu, Y., and Zhang, J.
\newblock Diffusion models for reinforcement learning: A survey.
\newblock \emph{arXiv:2311.01223}, 2023.

\bibitem[Albergo et~al.(2023)]{albergo2023building}
Albergo, M.~S., Boffi, N.~M., and Vanden-Eijnden, E.
\newblock Building normalizing flows with stochastic interpolants.
\newblock \emph{ICLR}, 2023.

\bibitem[Anderson(1982)]{anderson1982reverse}
Anderson, B.~D.~O.
\newblock Reverse-time diffusion equation models.
\newblock \emph{Stochastic Processes and their Applications}, 12(3):313--326, 1982.

\bibitem[Berner et~al.(2022)]{berner2022optimal}
Berner, J., Richter, L., and Lorenz, K.
\newblock An optimal control perspective on diffusion-based generative modeling.
\newblock \emph{TMLR}, 2022.

\bibitem[Black et~al.(2024)]{black2024training}
Black, K., Janner, M., Du, Y., Kostrikov, I., and Levine, S.
\newblock Training diffusion models with reinforcement learning.
\newblock \emph{ICLR}, 2024.

\bibitem[Chi et~al.(2023)]{chi2023diffusion}
Chi, C., Feng, S., Du, Y., Xu, Z., Cousineau, E., Burchfiel, B., and Song, S.
\newblock Diffusion policy: Visuomotor policy learning via action diffusion.
\newblock \emph{RSS}, 2023.

\bibitem[Davis et~al.(2025)]{davis2025gfm}
Davis, O., Albergo, M.~S., Boffi, N.~M., Bronstein, M.~M., and Bose, A.~J.
\newblock Generalised flow maps for few-step generative modelling on {R}iemannian manifolds.
\newblock \emph{arXiv:2510.21608}, 2025.

\bibitem[De~Bortoli et~al.(2022)]{debortoli2022riemannian}
De~Bortoli, V., Mathieu, E., Hutchinson, M., Thornton, J., Teh, Y.~W., and Doucet, A.
\newblock Riemannian score-based generative modelling.
\newblock \emph{NeurIPS}, 2022.

\bibitem[Deng et~al.(2026)]{deng2026drifting}
Deng, M., Li, H., Li, T., Du, Y., and He, K.
\newblock Generative modeling via drifting.
\newblock \emph{arXiv:2602.04770}, 2026.

\bibitem[Dhariwal and Nichol(2021)]{dhariwal2021diffusion}
Dhariwal, P. and Nichol, A.
\newblock Diffusion models beat {GAN}s on image synthesis.
\newblock \emph{NeurIPS}, 2021.

\bibitem[Fan et~al.(2024)]{fan2024reinforcement}
Fan, Y., Watkins, O., Du, Y., Liu, H., Ryu, M., Boutilier, C., Abbeel, P., Ghavamzadeh, M., Lee, K., and Lee, K.
\newblock Reinforcement learning for fine-tuning text-to-image diffusion models.
\newblock \emph{NeurIPS}, 2024.

\bibitem[Geng et~al.(2025)]{geng2025meanflow}
Geng, Z., Deng, M., Bai, X., Kolter, J.~Z., and He, K.
\newblock Mean flows for one-step generative modeling.
\newblock \emph{arXiv:2505.13447}, 2025.

\bibitem[Ho et~al.(2020)]{ho2020denoising}
Ho, J., Jain, A., and Abbeel, P.
\newblock Denoising diffusion probabilistic models.
\newblock \emph{NeurIPS}, 2020.

\bibitem[Janner et~al.(2022)]{janner2022planning}
Janner, M., Du, Y., Tenenbaum, J.~B., and Levine, S.
\newblock Planning with diffusion for flexible behavior synthesis.
\newblock \emph{ICML}, 2022.

\bibitem[Lipman et~al.(2023)]{lipman2023flow}
Lipman, Y., Chen, R.~T.~Q., Ben-Hamu, H., Nickel, M., and Le, M.
\newblock Flow matching for generative modeling.
\newblock \emph{ICLR}, 2023.

\bibitem[Liu et~al.(2023)]{liu2023flow}
Liu, X., Gong, C., and Liu, Q.
\newblock Flow straight and fast: Learning to generate and transfer data with rectified flow.
\newblock \emph{ICLR}, 2023.

\bibitem[Mishra et~al.(2026)]{mishra2026cdgs}
Mishra, U.~A., He, D., Chen, Y., and Xu, D.
\newblock Compositional diffusion with guided search for long-horizon planning.
\newblock \emph{arXiv:2601.00126}, 2026.

\bibitem[Potaptchik et~al.(2026)]{potaptchik2026meta}
Potaptchik, P., Saravanan, A., Mammadov, A., Prat, A., Albergo, M.~S., and Teh, Y.~W.
\newblock Meta flow maps enable scalable reward alignment.
\newblock \emph{arXiv:2601.14430}, 2026.

\bibitem[Shao et~al.(2024)]{shao2024deepseekmath}
Shao, Z., Wang, P., Zhu, Q., Xu, R., Song, J., Zhang, M., Li, Y., Wu, Y., and Guo, D.
\newblock {DeepSeekMath}: Pushing the limits of mathematical reasoning in open language models.
\newblock \emph{arXiv:2402.03300}, 2024.

\bibitem[Singhal et~al.(2025)]{singhal2025fksteering}
Singhal, R., Horvitz, Z., Teehan, R., Ren, M., Yu, Z., McKeown, K., and Ranganath, R.
\newblock A general framework for inference-time scaling and steering of diffusion models.
\newblock \emph{arXiv:2501.06848}, 2025.

\bibitem[Snell et~al.(2024)]{snell2024scaling}
Snell, C., Lee, J., Xu, K., and Kumar, A.
\newblock Scaling {LLM} test-time compute optimally can be more effective than scaling model parameters.
\newblock \emph{arXiv:2408.03314}, 2024.

\bibitem[Song et~al.(2021a)]{song2021scorebased}
Song, Y., Sohl-Dickstein, J., Kingma, D.~P., Kumar, A., Ermon, S., and Poole, B.
\newblock Score-based generative modeling through stochastic differential equations.
\newblock \emph{ICLR}, 2021.

\bibitem[Song et~al.(2021b)]{song2021sde}
Song, Y., Sohl-Dickstein, J., Kingma, D.~P., Kumar, A., Ermon, S., and Poole, B.
\newblock Score-based generative modeling through stochastic differential equations.
\newblock \emph{ICLR}, 2021.

\bibitem[Song et~al.(2023)]{song2023consistency}
Song, Y., Dhariwal, P., Chen, M., and Sutskever, I.
\newblock Consistency models.
\newblock \emph{ICML}, 2023.

\bibitem[Uehara et~al.(2024)]{uehara2024finetuning}
Uehara, M., Zhao, Y., Black, K., Hazan, E., Lee, J.~D., Jiang, N., Wang, Z., Levine, S., and Bietti, A.
\newblock Fine-tuning of continuous-time diffusion models as entropy-regularized control.
\newblock \emph{arXiv:2402.15194}, 2024.

\bibitem[Vincent(2011)]{vincent2011connection}
Vincent, P.
\newblock A connection between score matching and denoising autoencoders.
\newblock \emph{Neural Computation}, 23(7):1661--1674, 2011.

\bibitem[Zheng et~al.(2025)]{zheng2025diffusionnft}
Zheng, K., Chen, H., Ye, H., Wang, H., Zhang, Q., Jiang, K., Su, H., Ermon, S., Zhu, J., and Liu, M.-Y.
\newblock {DiffusionNFT}: Online diffusion reinforcement with forward process.
\newblock \emph{arXiv:2509.16117}, 2025.

\bibitem[Zhou et~al.(2025)]{zhou2025tvm}
Zhou, L., Parger, M., Haque, A., and Song, J.
\newblock Terminal velocity matching.
\newblock \emph{arXiv:2511.19797}, 2025.

\end{thebibliography}

\end{document}
